% ---------------------------------------------------------------------
% STATUT : RÉDIGÉ. Bloc de rigueur formelle destiné au jury quantitatif.
% Les trois propositions sont démontrées ET vérifiées numériquement par
% exploratory/vasicek_lab/6_contagion/32_proprietes_formelles.py (figure Z3).
% Requiert l'environnement `cle` (voir macros_requises.tex).
% ---------------------------------------------------------------------
\chapter{Trois propriétés du modèle}
\label{ch:proprietes}

\begin{cle}
\textbf{Objet du chapitre.} Le mémoire ne prétend pas apporter un théorème nouveau
d'analyse stochastique. Il établit en revanche trois propriétés de son propre objet, et
chacune est accompagnée de sa \textbf{vérification numérique} : l'énoncé dit ce qui doit
se produire, le calcul montre que cela se produit. La troisième est la plus importante,
car elle transforme la limite d'identification en résultat opposable.
\end{cle}

\section{La criticité dépend de l'ordre, non de l'ensemble}

Notons $\mathfrak{S}_n$ l'ensemble des permutations de $n$ éléments et, pour une chaîne
ordonnée $\sigma=(i_1,\dots,i_n)$ de piliers deux à deux distincts, $\mathcal{E}(\sigma)$
l'\emph{ensemble} $\{i_1,\dots,i_n\}$ de ses éléments. La vraisemblance de la chaîne s'écrit
\begin{equation}
  \mathcal{L}(\sigma) \;=\; \mathrm{ROOT}[i_1]\;\prod_{k=1}^{n-1}\mathrm{TRANS}[i_k][i_{k+1}] ,
  \label{eq:vraisemblance-chaine}
\end{equation}
et la criticité $c(\sigma)$ en est une fonction croissante, composée le long de l'ordre.

\begin{proposition}[Dépendance à l'ordre]\label{prop:ordre}
La criticité n'est pas une fonction de l'ensemble des piliers touchés : il existe
$\sigma$ et $\sigma'$ avec $\mathcal{E}(\sigma)=\mathcal{E}(\sigma')$ tels que
$c(\sigma)\neq c(\sigma')$. En particulier, $c$ n'est pas invariante par permutation.
\end{proposition}

\begin{proof}
Il suffit d'un contre-exemple, et la construction en fournit une famille. Dans
\eqref{eq:vraisemblance-chaine}, le produit $\prod_k \mathrm{TRANS}[i_k][i_{k+1}]$ dépend de
l'ordre dès que $\mathrm{TRANS}$ n'est pas symétrique, puisque permuter deux indices consécutifs
$i_k, i_{k+1}$ remplace le facteur $\mathrm{TRANS}[i_k][i_{k+1}]$ par
$\mathrm{TRANS}[i_{k+1}][i_k]$. Pour la paire $\{P_1,P_2\}$, le classeur donne
$\mathrm{TRANS}[1][2]=0{,}80$ et $\mathrm{TRANS}[2][1]=0{,}20$, d'où des criticités de $8$ et
$2$ pour le même ensemble $\{P_1,P_2\}$.
\end{proof}

\begin{cle}
\textbf{La contraposée est l'argument utile.} Toute mesure de risque $m$ qui ne dépend que de
l'ensemble touché vérifie $m(\sigma)=m(\sigma\circ\pi)$ pour tout $\pi\in\mathfrak{S}_n$. La
proposition~\ref{prop:ordre} exhibe donc un objet que \emph{aucune} telle mesure ne reproduit,
et l'hypothèse requise est minimale : il suffit que $\mathrm{TRANS}$ soit asymétrique, sans
qu'aucune \emph{valeur} de la direction ne soit revendiquée.
\end{cle}

\paragraph{Vérification.} Sur les $26$ sous-ensembles de taille au moins deux,
\textbf{$22$ changent de criticité selon l'ordre}, avec un écart atteignant \textbf{$6$
points sur une échelle de $10$}, soit $60\,\%$ de l'échelle, pour un \emph{même} ensemble
de piliers.

\paragraph{Portée.} Toute dépendance symétrique (copule) et tout score additif sont, par
construction, \textbf{invariants par permutation}. Aucun ne peut donc reproduire la
criticité. C'est l'argument d'exclusion des approches concurrentes, et il ne requiert
\emph{aucun} jugement d'expert sur la \emph{valeur} de la direction : il suffit que
$\mathrm{TRANS}$ ne soit pas symétrique.

\begin{cle}
\textbf{Ce qui n'est pas revendiqué, et pourquoi le dire.} Une propriété plus forte, la
\textbf{non-transitivité} de la dominance entre piliers, a été testée dans trois
formulations : la transitivité des corrélations qu'impose un modèle à facteur unique, la
positivité de la dépendance symétrisée, et l'existence de cycles de dominance par paire.
\textbf{Aucune ne se vérifie} : zéro violation de l'inégalité de transitivité sur $420$
triplets, valeur propre minimale $+0{,}405$ donc matrice de corrélation parfaitement
valide, et zéro cycle. Cette revendication est donc \textbf{retirée}. La dépendance à
l'ordre exclut exactement les mêmes concurrents et, elle, se démontre.
\end{cle}

\section{La normalisation de Leontief borne la propagation}

\begin{proposition}[Borne de propagation]\label{prop:leontief}
Pour une cascade indépendante de matrice moyenne $W\geq 0$, où $W_{jk}$ est la probabilité que
le pilier $j$ contamine directement $k$, et en notant $S_j$ l'ensemble des piliers atteints
depuis l'amorce $j$ :
\begin{equation}
  \E\bigl[\,|S_j|\,\bigr] \;\leq\; \bigl[(I-W)^{-1}\mathbf{1}\bigr]_j ,
  \label{eq:borne-leontief}
\end{equation}
le membre de droite étant fini si et seulement si $\rho(W) < 1$.
\end{proposition}

\begin{proof}
Le nombre moyen de descendants de génération $k$ issus de $j$ est la $j$-ème ligne de $W^{k}$.
La descendance totale espérée, comptée \emph{avec} multiplicité, vaut donc
$\sum_{k\geq 0} W^{k} = (I-W)^{-1}$, la série géométrique matricielle convergeant si et
seulement si $\rho(W)<1$. L'ensemble atteint compte chaque pilier une seule fois, là où la
descendance le compte autant de fois qu'il est contaminé, d'où la domination de $|S_j|$ par la
descendance totale.
\end{proof}

\begin{proposition}[Sous-criticité par construction]\label{prop:souscritique}
Soit $\mathrm{TRANS}\geq 0$, $s_j=\sum_k \mathrm{TRANS}[j][k]$ la somme de sa $j$-ème ligne, et
\begin{equation}
  W \;=\; g\,\frac{\mathrm{TRANS}}{\max_k s_k},
  \qquad g\in(0,1).
  \label{eq:normalisation-leontief}
\end{equation}
Alors $\rho(W)\leq g<1$, et la conclusion de la proposition~\ref{prop:leontief} vaut sans
condition supplémentaire.
\end{proposition}

\begin{proof}
Pour une matrice à coefficients positifs, le rayon spectral est majoré par toute norme
matricielle subordonnée, en particulier par la norme infinie, c'est-à-dire le maximum des sommes
de lignes : $\rho(W)\leq \max_j \sum_k W_{jk}$. Or, par \eqref{eq:normalisation-leontief},
$\sum_k W_{jk} = g\,s_j/\max_k s_k \leq g$ pour tout $j$, avec égalité pour la ligne réalisant
le maximum. D'où $\rho(W)\leq g<1$.
\end{proof}

\paragraph{Portée.} La sous-criticité n'est donc pas le fruit d'une calibration heureuse mais une
\textbf{conséquence de la normalisation} : c'est ce qui autorise la forme réduite $(I-W)^{-1}$
malgré les cycles de propagation, là où un graphe bayésien exigerait l'acyclicité. C'est aussi ce
qui sépare cette construction d'une matrice posée au hasard, pour laquelle rien ne garantirait
l'extinction de la cascade.

\paragraph{Vérification, et ce que l'écart mesure.} Avec le classeur actuel,
$\rho(W) = 0{,}506$. L'inégalité est vérifiée sur les cinq amorces et sur $300$ matrices
de l'ensemble admissible, sans aucun contre-exemple. L'écart entre la borne et la valeur
exacte va de $10{,}6\,\%$ (amorce P5) à $19{,}9\,\%$ (amorce P1). \textbf{Cet écart n'est
pas une imprécision} : il mesure exactement les \emph{collisions}, c'est-à-dire les
piliers atteints par plusieurs chemins, que la forme linéaire compte plusieurs fois.
Sur cinq piliers fortement connectés il n'est pas négligeable, ce qui justifie de
conserver le calcul exact par sous-ensembles plutôt que la seule forme close.

\section{La direction n'est pas identifiée, et cela coûte}

\begin{proposition}[Décomposition, et indistinguabilité observationnelle]\label{prop:decomposition}
Toute matrice réelle $W$ s'écrit de façon \emph{unique} $W = S + A$ avec
$S=\tfrac12(W+W^{\top})$ symétrique et $A=\tfrac12(W-W^{\top})$ antisymétrique, et ces deux
parts sont orthogonales pour le produit scalaire de Frobenius :
\begin{equation}
  \langle S, A\rangle_F \;=\; \mathrm{tr}(S^{\top}A) \;=\; 0 .
  \label{eq:orthogonalite}
\end{equation}
En outre $W$ et $W^{\top}$ ont la même partie symétrique et des parties antisymétriques
opposées. Par conséquent, toute statistique $T$ ne dépendant de $W$ que par $S$ vérifie
$T(W)=T(W^{\top})$ : elle ne peut pas distinguer $W$ de sa transposée.
\end{proposition}

\begin{proof}
L'existence est donnée par les formules ci-dessus. Pour l'unicité, si $S+A = S'+A'$ avec $S,S'$
symétriques et $A,A'$ antisymétriques, alors $S-S' = A'-A$ est à la fois symétrique et
antisymétrique, donc nulle. Pour \eqref{eq:orthogonalite}, la symétrie de $S$ et
l'antisymétrie de $A$ donnent
$\mathrm{tr}(S^{\top}A)=\mathrm{tr}(SA)=\mathrm{tr}\bigl((SA)^{\top}\bigr)
=\mathrm{tr}(A^{\top}S^{\top})=-\mathrm{tr}(AS)=-\mathrm{tr}(SA)$, d'où
$\mathrm{tr}(SA)=0$. Enfin $W^{\top}=S-A$, ce qui établit la dernière assertion.
\end{proof}

\paragraph{Conséquence.} La co-occurrence observée entre piliers est précisément une telle
statistique, fonction de $S$ seul. Le test directionnel ne rejette d'ailleurs pas l'hypothèse
$A = 0$
(chapitre~\ref{ch:identifiabilite} : asymétrie observée $119$ contre un nul de
permutation à $128 \pm 27$, soit $z = -0{,}33$, et aucune arête significative sur $18$).
Les deux matrices sont donc \textbf{également compatibles avec ce que l'on observe}. Elles
ne produisent pourtant ni le même capital, ni la même décision.

\paragraph{Vérification.} Retourner la matrice déplace le capital de $3{,}6\,\%$ et la perte
moyenne de $7{,}9\,\%$ (les niveaux sont établis en partie~\ref{part:resultats} ; seuls les
écarts importent ici). Surtout, \textbf{la priorité de remédiation bascule de P1 à P4}. Sur
l'ensemble admissible, l'écart entre une matrice et sa transposée atteint
$1\,173$~M\euro{}.

\begin{cle}
\textbf{L'énoncé à retenir : observationnellement indistinguables, matériellement
différentes.} C'est la formulation la plus nette de la frontière d'identification, et
elle ne repose sur \emph{aucun} dire d'expert. Elle chiffre en euros de capital, et en
inversion de priorité, le prix de ne pas savoir dans quel sens la contagion circule.
C'est ce qui fonde à la fois le modèle en couches (chapitre~\ref{chap:partielle}) et la
spécification de reporting adressée au régulateur.
\end{cle}

% =====================================================================
